\scanpage{64}
% archive foliation: 63
contient pas \uline{assez} de \inserted{$\Sigma$-}projectifs: il faut y rajouter $\Psi^0$\ldots) \deleted{Il est vrai que} Ces assertions sont d'ailleurs vraies aussi dans la plus gde catégorie $\unsure{\mathcal P|P}$ [détermination des \inserted{$\Sigma$-}proj et des $\Sigma$-inj, et le fait qu'il y en a assez].

Les $\operatorname{Ext}^i$ calculés par les résolutions injectives \inserted{à droite} ou projectives \inserted{à g.} \deleted{\unsure{sont} (à gauche} \unsure{ou les} deux simultanément, sont les \unsure{mêmes} pour $i\geq1$, ce sont les hyperext habituels de \inserted{$L_\bullet$ et de $M_\bullet$}, (pour $i=0$ on doit prendre les $\unsure{L^i|L^\bullet}$.

Si $F:\ Q=\unsure{\mathcal P|P}_1\subset \unsure{\mathcal P|P}\ \unsure{\longrightarrow|{-}}\ B$ (\unsure{$k$-addit.}) \deleted{est un} foncteur, on peut construire les foncteurs dérivés \inserted{à droite et à g.} (en utilisant des résolutions injectives ou projectives.

Nous sommes intéressés notamment aux $R^iF(\Psi^0)$, qu'on note aussi $\unsure{\mathcal H^i(F)|H^i(F)}$ [on doit pouvoir montrer que ce sont les dérivés \inserted{droits} de $F\mapsto F(\Psi^0)=\Gamma(F)$, quand $F$ varie dans la cat. des \deleted{foncteurs} \inserted{$k$-add.} $\unsure{\mathcal P|P}\unsure{\longrightarrow|{-}}B$ \inserted{$\Sigma$-}exacts à gauche; (Si on se restreint à de tels $F$); \unsure{qu’elle|en effet} cette catégorie est équivalente à \deleted{\unsure{celle}} \deleted{$\unsure{B}$} $\unsure{B^\bullet|B'}$, $F\mapsto F(\Psi^0)$ s'identifie \deleted{alors} au foncteur $C^\bullet\mapsto Z^0(C^\bullet)$\inserted{$=H^0(C^\bullet)$} sur $\unsure{B^\bullet|B'}$, dont les dérivés sauf erreur sont les $H^i(C^\bullet)$ pour $i\geq1$. On a aussi $RF(\Psi^0)\simeq F(C^\bullet)$

